\section{From Local Touch to Cross-Interface Representations}
\label{sec:representations}

\subsection{A hierarchy of representation scales}

Tactile models can be placed on a hierarchy from sensor-local appearance to
task-level physical relations (
\cref{fig:representation_hierarchy}). Progress at a lower level is enabling but
does not imply progress at every higher level. For example, a sensor-invariant
embedding may improve local contact classification while discarding absolute
orientation or force magnitude that a load-transfer model needs. Conversely, a
task-specific relational policy may model useful interactions while inheriting a
fragile local encoder.

\begin{figure}[t]
\centering
\begin{tikzpicture}[
  level/.style={draw=reviewblue, rounded corners, fill=reviewlight,
                minimum height=11mm, align=center, font=\small},
  goal/.style={draw=revieworange, rounded corners, fill=reviewwarm,
               minimum height=11mm, align=center, font=\small},
  arr/.style={-{Latex[length=2mm]}, thick, draw=reviewgray},
  node distance=4mm
]
\node[level, minimum width=25mm] (raw) {raw taxels /\\tactile frames};
\node[level, minimum width=25mm, right=of raw] (local) {local patch\\embedding};
\node[level, minimum width=25mm, right=of local] (body) {finger/link/hand\\tokens};
\node[level, minimum width=25mm, right=of body] (entity) {entity-anchored\\interfaces};
\node[goal, minimum width=27mm, right=of entity] (relation) {dynamic relations\\and load paths};
\draw[arr] (raw)--(local);
\draw[arr] (local)--(body);
\draw[arr] (body)--(entity);
\draw[arr] (entity)--(relation);

\node[font=\scriptsize, align=center, below=3mm of raw] {denoise, mask,\\calibrate};
\node[font=\scriptsize, align=center, below=3mm of local] {geometry, force,\\slip, semantics};
\node[font=\scriptsize, align=center, below=3mm of body] {kinematics and\\temporal context};
\node[font=\scriptsize, align=center, below=3mm of entity] {object/task frame,\\contact identity};
\node[font=\scriptsize, align=center, below=3mm of relation] {phase, load share,\\future wrench, risk};
\end{tikzpicture}
\caption{Representation hierarchy. The arrows denote additional organization,
not a mandatory neural architecture. Cross-interface claims require evidence at
the rightmost levels, even when the local encoder is strong.}
\label{fig:representation_hierarchy}
\end{figure}

\subsection{Local self-supervised and transferable touch}

Self-supervised tactile learning replaces expensive per-task labels with
objectives such as reconstruction, masked prediction, temporal consistency, and
instance discrimination. T-Dex couples pretraining to robotic play, showing that
unlabeled interaction can benefit downstream dexterity~\cite{guzey2023tdex}.
Sparsh scales self-supervision over vision-based tactile data and evaluates a
family of downstream tasks~\cite{higuera2025sparsh}. These approaches establish
that generic visual backbones are not enough and that tactile-specific pretraining
can be reused.

Transfer across sensors is harder because pixels encode different elastomers,
illumination, marker patterns, optics, and spatial scales. UniTouch learns unified
multimodal tactile representations~\cite{yang2024unitouch}; AnyTouch explicitly
addresses static and dynamic information across multiple visuo-tactile sensors
~\cite{feng2025anytouch}; T3 studies transferable Transformers
~\cite{zhao2024t3}; and UniT uses a unified tactile modelling approach
~\cite{xu2024unit}. Sparsh-X broadens the multisensory alignment setting
~\cite{higuera2025sparshx}, while TactAlign maps human tactile-glove features into
a robot tactile space for cross-embodiment transfer~\cite{tactalign2026}. The
common scientific question is not simply whether
embeddings transfer, but which tactile factors remain identifiable after
alignment. A sensor-invariant latent that collapses shear direction or scale may
be detrimental to load-path reasoning.

Whole-hand and morphology-aware models extend the unit of representation. HTT
uses heterogeneous tactile tokenization and masked learning~\cite{bi2026htt};
FTP-1 incorporates sensor morphology in a large tactile policy
~\cite{yuan2026ftp1}; HT-Bench supplies a whole-hand benchmark with egocentric
vision~\cite{huang2026htbench}. These works motivate explicit metadata for sensor
location, normal direction, link identity, and quality. In bimanual systems,
morphology-awareness must also prevent a shortcut in which ``left'' and ``right''
sensor IDs stand in for task role.

\subsection{Temporal tactile representations}

Touch is an event stream. Static contact appearance can be ambiguous: the same
deformation may correspond to loading or unloading, stable stick or the moment
before slip, depending on recent history. TactileVAD highlights geometric
aliasing in learned tactile dynamics~\cite{oller2023tactilevad}. Reactive control
systems increasingly split slow semantic reasoning from fast tactile reaction;
T-Rex uses tactile-reactive control at different time scales~\cite{niu2026trex},
while recent slow--fast visuo-tactile diffusion policies pursue a related control
principle. ViTacFormer also includes future-touch prediction as part of
cross-modal representation learning~\cite{heng2025vitacformer}.

Temporal modelling choices include fixed windows, recurrent states, temporal
convolution, causal Transformers, and state-space models. Their comparison is
often confounded by receptive field, parameter count, and action history. For
contact-transition evaluation, causal alignment is essential: a future tactile
frame may be a valid training target, but it must not leak into the deployed
encoder input. Reporting should distinguish synchronous state estimation,
$100$--$300$ ms forecasting, and closed-loop reaction.

\subsection{Multimodal fusion and spatial anchoring}

Tactile measurements acquire meaning from proprioception and vision. Early or
late concatenation is a strong baseline because a sufficiently expressive
temporal model can learn cross-modal interactions. Cross-attention gives tactile
tokens selective access to image, state, or action tokens; ViTacFormer is a
prominent recent example~\cite{heng2025vitacformer}. The fair question is whether
a proposed structure improves over this baseline at equal observation, compute,
and temporal context.

Spatial anchoring transforms local touch into a shared frame. SaTA anchors tactile
awareness using forward kinematics, enabling spatial reasoning beyond sensor
index~\cite{huang2025sata}. Object-centric approaches similarly normalize motion
with respect to the manipulated entity~\cite{chen2024objectcentric}. Anchoring is
particularly valuable for two hands because corresponding tactile features lie
in separate kinematic trees. Yet anchoring assumes that the relevant transform is
known or perceived; errors in link calibration and object pose should be tested,
not hidden.

\subsection{What existing objectives supervise}

\begin{table}[t]
\centering
\caption{Representation families, what they organize, and the unresolved
cross-interface question. Entries are representative rather than exhaustive.}
\label{tab:representation_families}
\begin{tabularx}{\linewidth}{P{0.18\linewidth} P{0.23\linewidth} P{0.25\linewidth} Y}
\toprule
\textbf{Family} & \textbf{Representative works} & \textbf{Primary organization} & \textbf{Unresolved bimanual question} \\
\midrule
Local SSL & T-Dex, Sparsh~\cite{guzey2023tdex,higuera2025sparsh} & contact appearance, geometry, dynamics & Does the latent retain load magnitude/direction and interface identity? \\
Cross-sensor alignment & UniTouch, AnyTouch, T3, Sparsh-X, TactAlign~\cite{yang2024unitouch,feng2025anytouch,zhao2024t3,higuera2025sparshx,tactalign2026} & shared tactile factors across hardware & Which physically relevant factors are lost to invariance? \\
Whole-hand / morphology-aware & HTT, HT-Bench, FTP-1~\cite{bi2026htt,huang2026htbench,yuan2026ftp1} & sensor layout, link/body coverage & Can roles and topology transfer across two heterogeneous hands? \\
Temporal dynamics & TactileVAD, T-Rex~\cite{oller2023tactilevad,niu2026trex} & event history and action-conditioned change & Are contact transitions predicted causally, before control action? \\
Cross-modal attention & ViTacFormer~\cite{heng2025vitacformer} & vision--touch--state token interaction & Does explicit interface structure add value to matched attention? \\
Spatial anchoring & SaTA~\cite{huang2025sata} & tactile features in kinematic space & How robust is the representation to pose error and topology change? \\
\bottomrule
\end{tabularx}
\end{table}

Pretraining objectives privilege different information. Reconstruction preserves
pixel detail but may waste capacity on illumination or marker appearance. Masked
prediction encourages context use but can exploit spatial redundancy. Contrastive
alignment emphasizes shared content but can suppress modality-specific physical
signals. Future prediction encodes dynamics but may average multimodal outcomes.
Task reward discovers control-relevant features but can exploit spurious state or
demonstration regularities. No objective is universally correct; its sufficiency
must be judged against the intended quantities and topology.

For bimanual coordination, promising auxiliary targets include contact onset and
release, future interface occupancy, load-share ratio, future local wrench, slip
risk, and calibrated uncertainty. These are not merely extra labels. They test
whether the latent exposes a mechanism that a downstream policy could use. Where
labels are available only in simulation, privileged-to-observation distillation
can train an observable student, but evaluation must verify that the student does
not require privileged inputs at test time.

\subsection{Necessary baselines for representation claims}

A new structured encoder should be compared against more than raw concatenation.
At minimum, a matched study should include: proprioception-only; tactile-only;
concatenation with the same local encoder; cross-attention; a temporal Transformer
or recurrent model; spatial anchoring where poses are available; a static graph;
the proposed inferred structure; and an oracle structure used only as an upper
bound. Parameter count, input history, optimization budget, augmentations, and
pretraining data must be controlled or reported. Without these controls, an
improvement may come from capacity, temporal context, or privileged supervision
rather than representation structure.
